% =====================================
\section{Le modèle et sa dérivation}
% =====================================

\subsection{Modèle de Malthus}\label{sec_modele_Malthus}
Thomas Malthus est considéré comme le premier ayant créé un modèle mathématique sur la croissance démographique. Le modèle est simple, on considère une population homogène, c'est-à-dire que les individus sont identiques d'un point de vu démographique. On considère le nombre total d'individu à l'instant $t \in \bbR_+$ que l'on note $P(t)$. On suppose que la population est isolée et vit dans un environnement stable avec des ressources illimitées. Alors, les variations de la population vont dépendre seulement du taux de fertilité et du taux de mortalité. On note $\beta$ le taux de naissance par habitant et $\mu$ le taux de mortalité par habitant. On suppose ces deux quantités constantes.
\\
Soient $t\in \bbR_+$ et $h \in \bbR_+^*$. On observe la variation de la population sur $[t,t+h]$. On a 
\begin{equation*}
    P(t+h) = P(t) + \text{naissances} - \text{décès}
\end{equation*}
\begin{itemize}
    \item le nombre de naissance sur $[t,t+h]$ est
    \begin{equation*}
        \beta P(t) h
    \end{equation*}
    \item le nombre de décès sur $[t,t+h]$ est 
    \begin{equation*}
        \mu P(t) h
    \end{equation*}
\end{itemize}
donc en substituant dans l'équation, on obtient 
\begin{equation*}
    P(t+h) = P(t) + \beta P(t) h - \mu P(t) h
\end{equation*}
en divisant par $h$ et en faisant tendre $h$ vers $0$ on obtient \textcolor{astral}{\textit{l'équation de Malthus}}
\begin{equation*}
    P'(t) = (\beta - \mu) P(t)
\end{equation*}
On note $\alpha:=\beta-\mu$ le \textcolor{astral}{\textit{paramètre de Malthus}}. La solution de cette EDO est donnée par 
\begin{equation*}
    P(t) = P(0)\ex^{\alpha t}
\end{equation*}
\paragraph{Interprétation.} Si le nombre de naissance est grand devant le nombre de mort \textit{i.e.} $\beta \gg \mu$ alors la population va croître de manière exponentielle. Si le nombre de naissance est petit devant le nombre de mort \textit{i.e.} $\beta \ll \mu$ alors la population va finir par s'éteindre.
\\\\
Cependant, ce modèle à ses limites, en effet il n'y a que trois cas possibles : 
\begin{enumerate}
    \item la population s'éteint,
    \item la population est stationnaire,
    \item la population ne cesse d'augmenter et le nombre d'individu tend vers $+\infty$.
\end{enumerate}
Il peut toutefois être utilisé pour des projections à court terme. Son principal intérêt réside dans le fait que le paramètre de Malthus $\alpha$ synthétise l'impact combiné des taux de natalité et de mortalité, ce qui indique clairement la tendance générale de la population : croissance ou décroissance. La principale raison pour laquelle les populations naturelles ne peuvent croître exponentiellement indéfiniment est que leur croissance est régulée par divers facteurs, notamment les limitations des ressources alimentaires et du territoire disponible. Reconnaissant l'incapacité du modèle de Malthus à saisir la dynamique démographique à long terme, d'autres modèles ont été étudié.

\begin{remark}
    On peut également faire évoluer le modèle de Malthus en ajoutant une fonction représentant la migration en fonction du temps, l'étude est la même et fait appel à la formule de Duhamel.
\end{remark}

\subsection{Modèle de McKendrick-Lotka}
L'objectif de cette section est d'introduire le modèle de référence pour l'étude démographique. Ce sont McKendrick \cite{McKendrick1926} et Lotka \cite{Lotka1911} qui sont à l'origine du modèle que nous allons étudier, mais nous utiliserons plutôt les travaux de Iannelli, Martcheva et Milner \cite{Iannelli2005} pour dériver le modèle. Il s'agit dans un premier temps de choisir les bons paramètres de description, leur donner un sens mathématique puis établir le modèle. On fait le choix de considérer seulement des modèles continus, la plupart de ces modèles ont leur analogue en discret (comme par exemple le modèle de Malthus) et donnent des résultats similaires.

\subsubsection{Motivation de la structure en âge}
Comme nous l'avons vu précédemment, le modèle de Malthus repose sur une hypothèse trop simplificatrice, il suppose une croissance exponentielle illimitée et traite la population comme un ensemble homogène, sans tenir compte de sa structure interne. Or, dans la réalité, les taux de natalité et mortalité ne sont pas uniformes, ils varient fortement selon l'âge des individus.
L'âge est l'un des paramètres les plus fondamentaux pour décrire une population, car il conditionne directement 
\begin{itemize}
    \item les capacités reproductives,
    \item les probabilités de survie,
    \item le comportement des individus dans la population.
\end{itemize}
Ignorer cette structure empêche de traiter des problèmes essentiels comme la planification de l'éducation, la gestion de la sécurité sociale ou les politiques de santé publique. On a quelques exemples de pyramides des âges dans différents pays dans \cite{Iannelli2017}. L'objectif est donc d'établir un modèle qui modélise l'évolution de la densité de population par âge au cours du temps, en incorporant une fécondité et une mortalité spécifiques à l'âge.

\subsubsection{Les éléments de description} 
On veut donc construire un modèle continu qui dépend de l'âge et du temps, pour ça on reprend les notations de \cite{Iannelli2017}. On va décrire l'évolution de la population par une fonction donnant la \textcolor{astral}{\textit{densité d'individu d'âge $a$ à l'instant $t$}}
\begin{equation*}
    x(a,t) \geq 0, \ \ \ \ \ a \in [0,a_{\max}], \ t \geq 0
\end{equation*}
où $a_{\max}$ est l'âge maximal de la population, supposé fini. On pourra considérer plus tard l'éventualité $a_{\max} = +\infty$. On note 
\begin{equation*}
    P(t) := \int_0^{a_{\max}} x(a,t)\d a
\end{equation*}
le nombre total d'individu à l'instant $t$. 
\\
On introduit la fonction $\beta$ qui représente le \textcolor{astral}{\textit{taux de fertilité de la population}} en fonction de l'âge. On peut alors considérer le nombre total de nouveau né à l'instant $t$
\begin{equation*}
    B(t):= \int_{0}^{a_{\max}}\beta(a)x(a,t)\d a
\end{equation*}
On introduit à présent la fonction $\mu$ qui représente le \textcolor{astral}{\textit{taux de mortalité}} de la population en fonction de l'âge. On peut alors considérer le nombre total de décès à l'instant $t$
\begin{equation*}
    D(t) := \int_{0}^{a_{\max}} \mu(a)x(a,t)\d a
\end{equation*}

\begin{remark}
    Ici, les fonctions $\beta$ et $\mu$ sont définies seulement 
    en fonction de l'âge, mais on pourrait les définir en fonction de l'âge et du temps afin d'envisager une évolution des taux de naissance et de mortalité dans le temps. Par exemple, si on étudie l'évolution de la population européenne de $1900$ à aujourd'hui il faudrait prendre en compte le fait que $\mu$ était différents en $1900$, en $1940$ et aujourd'hui.
\end{remark}

On définit ensuite la \textcolor{astral}{\textit{probabilité de survie}}, c'est-à-dire la probabilité pour qu'un nouveau né survive jusqu'à l'âge $a$
\begin{equation*}
    \Pi (a) := \exp\bigg(-\int_{0}^a \mu(\sigma)\d \sigma\bigg)
\end{equation*}
et on pose $\Pi(a_{\max}) = 0$ car on veut que pour un âge suffisament grand, ici $a_{\max}$, un individu n'ait aucune chance de survie.
\\
On définit de plus la \textcolor{astral}{\textit{fonction de maternité}} pour tout $a \in [0,a_{\max}]$
\begin{equation*}
    K(a) := \beta(a)\Pi(a)
\end{equation*}
elle synthétise la dynamique des naissances et des décès de la population.
\\
Enfin, on définit le \textcolor{astral}{\textit{nombre de reproduction}}
\begin{equation*}
    R := \int_0^{a_{\max}} \beta(a)\Pi(a)\d a = \int_0^{a_{\max}} K(a)\d a
\end{equation*}
il donne la moyenne du nombre de nouveau né qu'un individu produit au cours de sa vie. Ce paramètre jouera un rôle important dans le comportement asymptotique de la population. Intuitivement, si $R>1$ la population va augmenter et si $R<1$ la population va décroitre.

\begin{remark}
    On peut également définir le taux de mortalité $\mu$ et la probabilité de survie $\Pi$ à l'aide de probabilité comme dans \cite[p.~3]{Inaba2017}. En effet, soit $X$ la variable aléatoire représentant l'âge de mort d'un individu. On note $f$ la densité de $X$ et $F$ la fonction de répartition associée à $X$. Pour $a \in [0,a_{\max}]$, $\mu(a)$ représente le taux de mortalité de la population à l'âge $a$. Autrement dit, la probabilité qu'un individu d'âge $a$ décède dans un intervalle de temps infinitésimal $[a,a+h]$, sachant qu'il a déjà atteint l'âge $a$. On peut donc définir
    \begin{eqnarray*}
        \mu(a) &=& \lim_{h \to 0} \frac{\bbP(a\leq X \leq a+h \ | \ a \leq X)}{h} \\
        &=& \lim_{h \to 0} \frac{\bbP(a \leq X \leq a+h)}{\bbP(a \leq X)h} \\
        &=& \lim_{h \to 0} \frac{\bbP(a \leq X \leq a+h)}{\big(1-\bbP(X \geq a)\big)h}\\
        &=& \frac{f(a)}{1-F(a)} \\
        &=& - \frac{\d}{\d a}\ln{\big(1-F(a)\big)}
    \end{eqnarray*}
    On peut alors définir la probabilité de survie
    \begin{equation*}
        \Pi(a) = \bbP(a \leq X) = 1 - F(a)
    \end{equation*}
    Alors on a 
    \begin{equation*}
        \Pi'(a) = -f(a) = -\mu(a)\Pi(a)
    \end{equation*}
    Comme un individu d'âge $0$ est assuré de survivre, on a $\Pi(0)=1$ et donc on a le problème de Cauchy
    \begin{equation*}
        \left\{\begin{array}{l}
            \Pi'(a) = -\mu(a)\Pi(a) \\
            \Pi(0)=1
        \end{array}\right.
    \end{equation*}
    dont la solution est donnée par 
    \begin{equation*}
        \Pi(a) = \exp\bigg(-\int_0^a \mu(\sigma)\d \sigma\bigg)
    \end{equation*}
    En particulier, si $a_{\max}<+\infty$, on a $\Pi(a_{\max})=0$ et donc 
    \begin{equation*}
        \int_0^{a_{\max}} \mu(\sigma)\d \sigma = +\infty
    \end{equation*}
\end{remark}

\subsubsection{Modèle de McKendrick-Lotka}
On peut à présent définir le modèle de McKendrick-Lotka comme dans \cite{Iannelli2005}. Soient $t \in \bbR_+$, $a \in [0,a_{\max}]$ et $h>0$. Comme pour le modèle de Malthus, on observe la variation de la population entre $t$ et $t+h$. Cependant, la structure en âge apporte une configuration différente au modèle. Si une période de temps $h$ s'est écoulée, alors l'âge des individus a augmenté en passant de $a$ à $a+h$. D'où la variation de la population sur $[t,t+h]$
\begin{equation*}
    x(a+h,t+h) = x(a,t) - \text{décès}
\end{equation*}
Ici, comme on regarde la densité de la population et pas le nombre d'individu, seul les décès interviennent. En effet, on ne peut pas avoir la naissance d'individu d'âge $a$ directement. Le taux de décès entre $t$ et $t+h$ est donné par
\begin{equation*}
    \int_0^{h}\mu(a+s)x(a+s,t+s)\d s
\end{equation*}
En effectuant un changement de variable, on peut réécrire 
\begin{equation*}
    \int_t^{t+h}\mu(a+s-t)x(a+s-t,s)\d s
\end{equation*}
La variation de population devient 
\begin{equation*}
    x(a+h,t+h) = x(a,t) - \int_t^{t+h}\mu(a+s-t)x(a+s-t,s)\d s
\end{equation*}
On pose 
\begin{equation*}
    G(t) := \int_0^{t}\mu(a+s-t)x(a+s-t,s)\d s
\end{equation*}
d'où
\begin{equation*}
    x(a+h,t+h)-x(a,t+h)+\big(x(a,t+h) - x(a,t)\big) = - \big(G(t+h)-G(t)\big)
\end{equation*}
On divise par $h$ et on fait tendre $h$ vers $0$, on obtient alors 
\begin{equation*}
    \partial_ax(a,t)+\partial_tx(a,t) = -G'(t) = -\mu(a)x(a,t)
\end{equation*}

%Dans un premier temps, on considère la fonction 
%$N(a,t)$ définie par 
%\begin{equation*}
%    N(a,t) = \int_0^a x(\sigma,t)\d \sigma
%\end{equation*}
%Elle représente le nombre d'individu qui, à l'instant $t$, on 
%un âge entre $0$ et $a$. Soient $t \in \bbR_+$, $a \in 
%[0,a_{\max}]$ et $h>0$. Comme pour le modèle de Malthus, on 
%observe la variation de la population entre $t$ et $t+h$.
%Cependant, la structure en âge apporte une configuration différente
%au modèle. Si une période de temps $h$ s'est écoulée, alors 
%l'âge des individus à augmentés en passant de $a$ à 
%$a+h$. D'où la variation de la population sur 
%$[t,t+h]$
%\begin{equation*}
%    N(a+h,t+h) = N(a,t)+\text{naissances} - \text{décès}
%\end{equation*}
%On sait que le nombre total de naissance à l'instant $t$ est $B(t)$,
%donc le nombre de naissance entre $t$ et $t+h$ est donné par
%\begin{equation*}
%    \int_t^{t+h}B(s)\d s
%\end{equation*}
%On veut maintenant déterminer le nombre de décès, il faut déterminer
%comment l'âge varie sur la période $[t,t+h]$. On veut aller jusqu'à
%$a+h$ donc si $u \in [t,t+h]$ alors $u-t \in [0,h]$ et donc 
%l'âge $\sigma$ varie dans $[0,a+u-t]$. D'où le nombre total de 
%mort sur la période $[t,t+h]$
%\begin{equation*}
%    \int_t^{t+h}\int_0^{a+u-t}\mu(\sigma)x(\sigma,u)\d \sigma %\d u
%\end{equation*}
%Finalement, le bilan démographique devient 
%\begin{equation}\label{bilan_demo1}
%    N(a+h,t+h) = N(a,t) +\int_t^{t+h}B(s)\d s - \int_t^{t+h}\int_0^{a+u-t}\mu(\sigma)x(\sigma,u)\d \sigma \d u
%\end{equation}
%On pose 
%\begin{eqnarray*}
%    M(t) := \int_0^t B(s)\d s \ \ \ &\text{ et }& \ \ \ L(t) := \int_{0}^{t} \int_0^{a+u-t}\mu(\sigma)x(\sigma, u)\d \sigma \d u
%\end{eqnarray*}
%Alors \eqref{bilan_demo1} devient 
%\begin{equation*}
%    N(a+h,t+h)-N(a,t+h)+\big(N(a,t+h) - N(a,t)\big) = M(t+h)-M(t) + \big(L(t+h)-L(t)\big)
%\end{equation*}
%d'où en divisant par $h$ puis en passant à la limite quand $h$
%tend vers $0$ on obtient 
%\begin{eqnarray*}
%    \partial_aN(a,t)+\partial_tN(a,t) = M'(t)+L'(t)
%\end{eqnarray*}
%soit 
%\begin{equation*}
%    x(a,t) + \int_0^a \partial_t x(\sigma,t)\d \sigma = B(t) %- \int_0^a \mu(\sigma)x(\sigma,t)\d \sigma
%\end{equation*}
%enfin, on dérive l'équation par rapport à la variable $a$, 

Ce qui donne \textcolor{astral}{\textit{l'équation de McKendrick-Lotka}}
\begin{equation}\label{eq_McKendrick_Lotka}
    \partial_a x(a,t) + \partial_tx(a,t) = -\mu(a)x(a,t)
\end{equation}
En ajoutant une donnée initiale, représentant la nombre d'individu d'âge $a$ à l'instant $t=0$, $x_0(a)$ et la condition au bord donnant le nombre de nouveau né à l'instant $t \in \bbR_+^*$, on obtient le \textcolor{astral}{\textit{modèle de McKendrick-Lotka}}
\begin{equation}\label{modele_McKendrick_Lotka}
    \left\{\begin{array}{l}
        \partial_a x(a,t) + \partial_t x(a,t) = -\mu(a)x(a,t) \\
        \\
        x(0,t) = B(t) = \displaystyle{\int_0^{a_{\max}}} \beta (a)x(a,t)\d a \\
        \\
        x(a,0) = x_0(a)
    \end{array}\right.
\end{equation}

\begin{remark}
    En supposant $\mu$ et $\beta$ constante on peut retrouver le modèle de Malthus. En effet, il nous suffit d'intégrer \eqref{eq_McKendrick_Lotka} entre $0$ et $a_{\max}$
    \begin{equation*}
        \int_0^{a_{\max}} \partial_a x(a,t)\d a + \int_0^{a_{\max}} \partial_t x(a,t)\d a = - \mu \int_0^{a_{\max}} x(a,t)\d a
    \end{equation*}
    On permute intégrale et dérivée en temps et avec les notations introduites précédemment
    \begin{equation*}
        x(a_{\max},t)-x(0,t) + \frac{\d}{\d t} P(t) = -\mu P(t)
    \end{equation*}
    On a $x(a_{\max},t)=0$ et on utilise la condition initiale 
    \begin{eqnarray*}
        -\beta P(t) + P'(t) = -\mu P(t) &\Longrightarrow& 
        P'(t) = (\beta - \mu)P(t)
    \end{eqnarray*}
\end{remark}

\paragraph{Hypothèses.} Au vu du contexte démographique, il est raisonnable de faire les hypothèses suivantes 
\begin{enumerate}
    \item $\beta \in L^{\infty}(0,a_{\max})$ et pour presque tout $a \in [0,a_{\max}]$, $0 \leq \beta(a) \leq \beta_+$ ;
    \item comme $\Pi(a_{\max}) = 0$, on pose
    \begin{equation*}
        \int_0^{a_{\max}}\mu(\sigma)\d \sigma = +\infty
    \end{equation*}
    \item $\mu \in \Lloc{0,a_{\max}}$ et pour presque tout 
    $a \in [0,a_{\max}]$, $\mu(a) \geq 0$ ;
    \item $x_0 \in L^1(0,a_{\max})$ et pour presque tout $a \in [0,a_{\max}]$, $x_0(a) \geq 0$.
\end{enumerate}

\begin{remark}
    \begin{itemize}
        \item On pourrait donner un sens à un taux de natalité négatif en le définissant comme un taux de mortalité. On préfère cependant définir deux fonctions distinctes $\beta$ et $\mu$ que l'on suppose positive.
        \item On suppose que $\beta$ est essentiellement bornée pour contrôler le taux de natalité, il serait difficile de donner un sens démographique au fait que $\|\beta\|_{L^{\infty}} = +\infty$.
        \item $\mu$ dans $L^1_{\text{loc}}$ découle simplement du fait que l'âge maximal est supposé fini, donc la probabilité de survie jusqu'à l'age maximal, $\Pi(a_{\max})$, doit être supposée nulle pour assurer la cohérence avec la réalité, ce qui nous oblige, au vu de la définition de $\Pi$, de supposer
        \begin{equation*}
            \int_0^{a_{\max}}\mu(\sigma)\d \sigma = +\infty
        \end{equation*}
        et donc avoir $\mu$ dans $L^1_{\text{loc}}$.
        \item La densité de population initiale est supposée $L^1$ afin de pouvoir déterminer le nombre moyen d'individu à l'instant initial.
    \end{itemize}
\end{remark}